Prettig aangeschoten
maar niet afgeschoten,

bleef de jager rolvast
en hield haar stevig vast.

Emmy’s nood was niet gedoofd.
Het zwaard boven haar hoofd

hing aan haar paardenhaar.
Zijn greep was het gevaar

waardoor ze wankel was
bij het stille moeras.

Ze bleef opgehangen
sinds ze was gevangen.

Haar licht werd verduisterd.
Angst hield haar gekluisterd.

Deze kernemotie
stond nauw in connectie

met haar alerte brein,
dat zo niet vrij kon zijn.

Net als bij prooidieren,
die steeds hun vizieren

op elk gevaar richten
en zo spanning stichten.

Die hechte waakzaamheid,
die haar van binnen bijt,

brak elke golf van rust.
Geen context die het sust.

Een brandweerkazerne
in onze interne

geest die steeds actief is.
En is er eens niets mis

dan steekt een brandweerman
zelf de vlam in de pan.

Zo houdt hij zich bezig
en is zijn baan stevig.

Hij blijft vergelijken
zodat steeds zal blijken

dat iets goed of slecht is.
Groot of klein, raak of mis.

De relativiteit
vermomd als wettig feit.

Fenomenologen
hebben vooral ogen

voor een echte dreiging
met angst als waarneming.

Want ging het bij Emmy
wel om een illusie?

Emmy moest moed baren.
De weeën ervaren,

die haar hevige pijn
omzetten in wél-zijn.

Haar weemoed was een les.
Haar leed de keuzestress.

Jean-Paul Sartre stelde
dat vrijdom beknelde

indien bewust geduid.
Erich Fromm legde uit

in ‘De angst voor vrijheid’:
verantwoordelijkheid

roept radeloosheid op
en zoekt steun hogerop.

Doch Aristoteles
zag angst als wijze les.

Het vroeg moed om de vrees
voor haar geest versus vlees

onder ogen te zien.
Emmy kon dit inzien

ontwaakt uit de wilschok,
die zoetjesaan wegtrok.

Met adem als anker
werd haar huid niet blanker.

Ze wiegde haar lichaam
kalm, ritmisch maar langzaam.

“Mijn vrees heeft ook zijn pracht.
Immers de bijbel bracht:

Vrees God en volg zijn wet.
Wijsheid komt daar aan zet.

De vrees des Heren leert,
dat alle angst verteert.

Voor wie ontzag ervaart,
verdwijnt al wat bezwaart.”

Dit grootse “awe” gevoel
ervoer Emmy als koel.

Vrees was Emmy’s vijand,
ze wees het van de hand

en had het onderdrukt,
waardoor “awe” haar niet lukt.

Door het te negeren
kon ze niets meer leren

over de kern van angst.
De waarheid duurt het langst

en kwam plots opdagen
zonder eerst te vragen

of het dan wel schikte
en ze zich verslikte.

De spanning, ongekend,
verstoord en onbekend,

moest ze dan ongewend
op haar zwakste moment
plotseling mee omgaan.
Kort na Emmy’s bestaan

was dit nog een reflex.
Net als honger en sex.

Gezond en natuurlijk.
Het werd meer figuurlijk

naarmate ze leerde,
waakzaam observeerde,

gedwongen aanleerde
en communiceerde.

Ze erfde ook angsten
van de dwaze bangsten.

Ze kon vrees niet dragen.
Ze moest het verjagen

als een quasi-demon,
die ze uitdrijven kon.

Als een nare vijand,
die haar geest had bemand.

Ze was haar verzet moe:
“Als ik wederom doe

wat ik altijd al deed,
dan krijg ik ook het leed

wat ik altijd al kreeg”.
Hoewel zo haar moed steeg,

miste ze nog inzicht
en bleef ze een bang wicht.

“Papa bent u bevreesd
voor dat piepkleine beest

dat over de grond rent?”
Nadat hij had ontkend,

deed hij vlug sokken aan
alvorens te gaan staan.

“Emmy heb de moed pijn,
als die komt uit het brein,

met verstand te dragen.
Hoe kun je dus vragen

of een piepkleine mier
gelijk staat aan een stier?

Voor wie mijn bedeesdheid
juist getuigt van wijsheid.”

Moeder grinnikte zacht.
Ze dacht weer aan de nacht,

hartje zomer was het,
en zo heet in hun bed,

dat allen naakt sliepen.
Tot wat jongens riepen

“Brand, brand, kom naar buiten!”
Vuur kon je omsluiten,

in die brandbare tijd,
zo snel dat een doodsstrijd

al binnen seconden
bij mensen ontstonden.

Dus stond het gezin snel
ontkleed op het appèl.

Vader sloot later aan.
Hij trok eerst sokken aan.

Als een held op sokken
extra opgetrokken.

De kwajongens lachten.
Deze hoofdverdachten

van het vals alarm slaan,
bleven manhaftig staan.

Moeder uitte klachten.
“Ik moet altijd wachten.

Je bent nooit eens op tijd!”
“Als de dood ons ooit scheidt,

beloof ik dat ik zal wachten!”
Dat kon men verwachten

van vader als lafaard.
De pijn van een uitvaart

zo gehaaid te mijden.
Moeder zou niet lijden

onder vader zijn dood.
Wou juist dat hij opschoot.

"Fear Is a Man's Best Friend"
was als werk onbekend.

Angst voor het gapend gat
dat tussen denken zat.

Levensenergie drukt.
Prana werd onderdrukt.

Het brein wilde vluchten.
Ze had veel te duchten

in de begoocheling
waar haar brein naartoe ging.

Het ego was ontstaan
om bij angst weg te gaan.

“Daar heb ik geen last van!”,
was het ontkenningsplan.

In plaats van met verstand:
“Wat is er aan de hand?”

Vrees was een afknapper.
Schiet niet de boodschapper,

zoals de landheer deed.
De oorsprong van het leed

blijft dan geheimzinnig.
Het brein blijft onwennig

in het omgaan met schrik
en wordt een bangerik

in plaats van een leerling.
Vrees wees naar de richting.

Meestal kwam het gevaar
uiteindelijk uit haar.

Angst was het knalfestijn
van haar fantasierijk brein.

Alleen vierde Emmy
niet deze potentie.

De horror die ze schreef
en levenslang in bleef,

kon ze niet verrijken,
met humor naar kijken

of het script van schikken
om niet meer te schrikken.

Wordt haar angst in dit script
door Emmy aangestipt?
Gaf verdriet haar de kracht
om haar enge onmacht

toch wel te aanvaarden,
zodat ze kon aarden?
